\paragraph{可寻址 Gödel 外置的 Chebyshev 码长主项与见证位膨胀障碍}
在逐坐标可寻址的 Gödel 整数化已经获得精确乘积下界、平均复杂度最优律与最坏冗余发散之后，还可把最坏情形码长进一步精化为由 Chebyshev 函数控制的显式主项。

\begin{theorem}[可寻址 Gödel 外置的最坏码长下界由 Chebyshev 主项精确控制]\label{thm:xi-addressable-godel-worstcase-chebyshev-sharp-lower-bound}
\leanverified{paper\_xi\_addressable\_godel\_worstcase\_chebyshev\_sharp\_lower\_bound}
沿用定理 \ref{thm:xi-addressable-godel-product-lower-bound-classification} 与定理 \ref{thm:xi-addressable-godel-worstcase-redundancy-diverges} 的记号。对任意满足逐坐标可寻址性的单射
\[
E_T:\Sigma^T\hookrightarrow \mathbb N_{\ge 1},
\qquad
|\Sigma|=q\ge 2,
\]
记
\[
L_T(E_T):=\max_{s\in\Sigma^T}\log E_T(s).
\]
若把定理 \ref{thm:xi-addressable-godel-product-lower-bound-classification} 中出现的互异素数按升序重排为
\[
\pi_1<\cdots<\pi_T,
\]
并记 \(\mathfrak p_t\) 为第 \(t\) 个素数，则
\[
L_T(E_T)
\ge
\sum_{t=1}^{T}\log \pi_t
\ge
\sum_{t=1}^{T}\log \mathfrak p_t
=
\vartheta(\mathfrak p_T).
\]
从而
\[
L_T(E_T)
\ge
T\log T+T\log\log T+O(T).
\]
\end{theorem}
\begin{proof}
沿用定理 \ref{thm:xi-addressable-godel-product-lower-bound-classification} 中的估值族记号
\[
\mathcal V_{t,a}:=
\{v_{p_t}(E_T(s)):\ s_t=a\},
\qquad
m_{t,a}:=\min \mathcal V_{t,a}.
\]
由于不同字母 \(a\neq b\) 的集合 \(\mathcal V_{t,a},\mathcal V_{t,b}\) 两两不交，故最小值 \(m_{t,a}\) 也彼此不同。特别地，在这 \(q\ge 2\) 个互异非负整数中，至多有一个等于 \(0\)。于是对每个 \(t\)，都可选取某个字母 \(a_t^\star\in\Sigma\) 使
\[
m_{t,a_t^\star}\ge 1.
\]
令
\[
s^\star:=(a_1^\star,\dots,a_T^\star)\in\Sigma^T.
\]
则对每个 \(t\) 都有
\[
v_{p_t}(E_T(s^\star))\ge m_{t,a_t^\star}\ge 1.
\]
因此
\[
L_T(E_T)
\ge
\log E_T(s^\star)
\ge
\sum_{t=1}^{T}\log p_t
=
\sum_{t=1}^{T}\log \pi_t.
\]

再由 \(\pi_1<\cdots<\pi_T\) 为 \(T\) 个互异素数，知其逐项满足
\[
\pi_t\ge \mathfrak p_t
\qquad (1\le t\le T),
\]
故
\[
\sum_{t=1}^{T}\log \pi_t
\ge
\sum_{t=1}^{T}\log \mathfrak p_t
=
\vartheta(\mathfrak p_T).
\]
这就得到第一条下界。

最后，由经典素数定理的精化形式，
\[
\mathfrak p_T
=
T\log T+T\log\log T+O(T),
\]
并且
\[
\vartheta(x)=x+o(x)\qquad (x\to\infty).
\]
代入 \(x=\mathfrak p_T\) 即得
\[
\vartheta(\mathfrak p_T)
=
T\log T+T\log\log T+O(T).
\]
于是所示精化下界成立。证毕。
\end{proof}

\leanverified{paper\_xi\_addressable\_witness\_no\_asymptotically\_linear\_bitlength}
\begin{corollary}[逐坐标 prime-addressable 见证外置不存在渐近线性位长度]\label{cor:xi-addressable-witness-no-asymptotically-linear-bitlength}
设 \(T=T(n)\to\infty\)，并对每个 \(n\) 给定满足逐坐标可寻址性的单射
\[
E_{T(n)}:\Sigma^{T(n)}\hookrightarrow \mathbb N_{\ge 1},
\qquad
|\Sigma|\ge 2.
\]
则
\[
\max_{s\in\Sigma^{T(n)}}\log_2 E_{T(n)}(s)
\ge
T(n)\log_2 T(n)+T(n)\log_2\log T(n)+O(T(n)).
\]
特别地，
\[
\max_{s\in\Sigma^{T(n)}}\log_2 E_{T(n)}(s)\neq O(T(n)).
\]
若进一步
\[
T(n)=n^c
\qquad (c>0),
\]
则更具体地有
\[
\max_{s\in\Sigma^{T(n)}}\log_2 E_{T(n)}(s)
\ge
c\,n^c\log_2 n+n^c\log_2\log n+O(n^c).
\]
\end{corollary}
\begin{proof}
把定理 \ref{thm:xi-addressable-godel-worstcase-chebyshev-sharp-lower-bound} 应用于 \(T=T(n)\)，并除以 \(\log 2\)，得到
\[
\max_{s\in\Sigma^{T(n)}}\log_2 E_{T(n)}(s)
\ge
\frac{1}{\log 2}
\Bigl(
T(n)\log T(n)+T(n)\log\log T(n)+O(T(n))
\Bigr),
\]
即首个不等式。由于右端严格超线性于 \(T(n)\)，故不可能属于 \(O(T(n))\)。

若 \(T(n)=n^c\)，则
\[
T(n)\log_2 T(n)=c\,n^c\log_2 n,
\]
且
\[
T(n)\log_2\log T(n)
=
n^c\log_2(c\log n)
=
n^c\log_2\log n+O(n^c).
\]
代回即可得到所示具体下界。证毕。
\end{proof}
